\subsection{Homogeneous extensions of representations: Problem 17.101}
\label{cand:17.101}
\problemstatement{17.101}

Fix a commutative ring \(k\) with identity. A representation
is a left \(kG\)-module \(V\) with its acting group \(G\);
an embedding is injective on both sorts. The relevant
homogeneity extends isomorphisms of subrepresentations
whose group is finitely generated and whose module is
finitely generated over that group's algebra. The coefficient
ring remains fixed; compare the category in \cite{gvaramia-plotkin}.

Given any set of partial isomorphisms
\((f_i,\theta_i):(U_i,A_i)\to(U_i',B_i)\) in \((V,G)\),
form the multiple HNN extension
\[
 H=\langle G,t_i:t_iat_i^{-1}=\theta_i(a),\
                         a\in A_i,\text{ all }i\rangle.
\]
The HNN embedding theorem embeds \(G\) in \(H\)
\cite{hnn}. Define
\[
 M=kH\otimes_{kG}V,\qquad
 E=\bigoplus_i kH\otimes_{kA_i}U_i,\qquad
 \delta_i(h\otimes u)=h\otimes u-ht_i^{-1}\otimes f_i(u),
\]
and \(W=M/\operatorname{im}\delta\).
The map \(\delta_i\) is well-defined because
\(at_i^{-1}=t_i^{-1}\theta_i(a)\) and \(f_i\) is equivariant.

We verify that no vector of \(V\) is lost. The HNN
Bass--Serre tree has vertices \(hG\) and edges \(hA_i\)
with endpoints \(hG,ht_i^{-1}G\). Its being a tree follows
from HNN normal form: a reduced closed path would be a
reduced stable-letter word lying in \(G\). As \(k\)-modules,
\(M\) is the direct sum of one copy of \(V\) at each
vertex, and \(E\) the direct sum of the edge modules.
Each edge maps injectively into both incident vertex
modules, using the original inclusion and \(f_i\).

If \(\delta(z)\) were supported at a single vertex \(v_0\)
for a nonzero \(z\in E\), the finitely many nonzero edge
components of \(z\) would form a finite forest with an edge.
It has a leaf other than \(v_0\). The single incident
nonzero edge contributes a nonzero vector at that leaf,
by injectivity, a contradiction. Hence
\(\operatorname{im}\delta\) intersects each vertex module
trivially, and \(V\hookrightarrow W\).
This argument uses no semisimplicity, purity, or
finite-dimensionality.

In \(W\), the edge relation gives \(t_i u=f_i(u)\).
The automorphism of \((W,H)\) with components
\[
 w\mapsto t_iw,\qquad h\mapsto t_iht_i^{-1}
\]
therefore extends the prescribed partial isomorphism.
Now repeat the construction at stages \(0,1,2,\ldots\),
each time including all isomorphisms between finitely
generated subrepresentations, and take the union.
The domain and range of any such isomorphism in the
union lie in a common stage: finitely many group and
module generators suffice. A stable letter at the next
stage realizes it, and its action continues to be an
automorphism of the union. This proves homogeneity.

If \(\kappa=\max(\aleph_0,|k|,|G|,|V|)\), there are at
most \(\kappa\) relevant finite tuples and maps at each
stage; induction, quotients, and the countable union
preserve that bound. Thus a countable input over a
countable ring has a countable homogeneous extension.
Source: \artifact{research/17.101-proof.md}.

